% arxiv-twocolumn.latex: the shared pandoc template for a two-column arXiv-style preprint.
% Built by build_paper.py (pandoc, then tectonic or any XeLaTeX); arxiv-twocolumn.lua is its filter.
% Engine: XeTeX (fontspec, unicode-math). Fonts: Latin Modern text and math, DejaVu Sans Mono for code;
% a character the current font lacks falls back to Latin Modern Math, then to DejaVu Sans.
\documentclass[10pt,letterpaper,twocolumn]{article}
\usepackage[letterpaper,top=0.75in,bottom=0.9in,left=0.7in,right=0.7in,columnsep=0.25in]{geometry}
\usepackage{amsmath}
\usepackage{fontspec}
\usepackage{unicode-math}
\setmainfont{lmroman10}[Extension=.otf,UprightFont=*-regular,BoldFont=*-bold,ItalicFont=*-italic,BoldItalicFont=*-bolditalic]
\setsansfont{lmsans10}[Extension=.otf,UprightFont=*-regular,BoldFont=*-bold,ItalicFont=*-oblique,BoldItalicFont=*-boldoblique]
\setmonofont{DejaVuSansMono}[Extension=.ttf,UprightFont=*,BoldFont=*-Bold,ItalicFont=*-Oblique,BoldItalicFont=*-BoldOblique,Scale=MatchLowercase]
\setmathfont{latinmodern-math.otf}
\newfontfamily\typesetfallbackmath{latinmodern-math.otf}
\newfontfamily\typesetfallbacksans{DejaVuSans}[Extension=.ttf,UprightFont=*,BoldFont=*-Bold,ItalicFont=*-Oblique,BoldItalicFont=*-BoldOblique]
% \typesetmap{"XXXX} (build_paper.py writes one for each character outside Latin-1 and Latin
% Extended-A that the paper uses) makes the character active at the start of the document: in text it
% comes from the current font when that has it, else from the first fallback that does; in math it
% keeps whatever unicode-math gave it (a prime stays a prime).
\makeatletter
\protected\def\typesetglyph#1{%
  \iffontchar\font#1\relax\char#1\relax
  \else{\typesetfallbackmath\iffontchar\font#1\relax\char#1\relax\else\typesetfallbacksans\char#1\relax\fi}%
  \fi}
\def\typeset@math#1{%
  \ifcsname typeset@m@#1\endcsname\csname typeset@m@#1\expandafter\endcsname\else\Uchar#1 \fi}
\protected\def\typesetchar#1{\ifmmode\expandafter\typeset@math\else\expandafter\typesetglyph\fi{#1}}
\def\typesetmap#1{%
  \edef\typeset@n{\number#1}% outside \lowercase, which would lowercase the hex digits A-F
  \begingroup\lccode`\~=\typeset@n\relax
  \lowercase{\endgroup
    \ifdefined~\expandafter\let\csname typeset@m@\typeset@n\endcsname~\fi
    \edef~{\noexpand\typesetchar{\typeset@n}}}%
  \catcode\typeset@n=13\relax}
\makeatother
% amssymb names that unicode-math spells differently.
\AtBeginDocument{%
  \providecommand{\square}{\mdlgwhtsquare}\providecommand{\Box}{\mdlgwhtsquare}%
  \providecommand{\blacksquare}{\mdlgblksquare}\providecommand{\lozenge}{\mdlgwhtlozenge}%
  \providecommand{\blacklozenge}{\mdlgblklozenge}}
% Legacy two-letter font commands (\rm, \bf, \it, ...), in text and in math.
\providecommand{\rm}{}\DeclareOldFontCommand{\rm}{\normalfont\rmfamily}{\mathrm}
\providecommand{\bf}{}\DeclareOldFontCommand{\bf}{\normalfont\bfseries}{\mathbf}
\providecommand{\it}{}\DeclareOldFontCommand{\it}{\normalfont\itshape}{\mathit}
\providecommand{\sf}{}\DeclareOldFontCommand{\sf}{\normalfont\sffamily}{\mathsf}
\providecommand{\tt}{}\DeclareOldFontCommand{\tt}{\normalfont\ttfamily}{\mathtt}
\usepackage{microtype}
\usepackage[dvipsnames]{xcolor}
\usepackage{graphicx}
\usepackage{booktabs}
\usepackage{array}
\usepackage{ragged2e}
\usepackage{xtab}
\usepackage{cuted}
\usepackage{caption}
\captionsetup{font=small,labelfont=bf,skip=4pt}
\captionsetup[table]{position=top}
\usepackage{fvextra}
\usepackage{xurl}
\usepackage{titlesec}
\usepackage{fancyhdr}
\usepackage{footmisc}
\usepackage{enumitem}
\setlist{itemsep=1pt,topsep=3pt,parsep=0pt}
\usepackage{hyperref}
\hypersetup{colorlinks=true,linkcolor=MidnightBlue,citecolor=MidnightBlue,urlcolor=MidnightBlue,filecolor=MidnightBlue,
  pdftitle={Selection for language-model context{,} read from the data},pdfauthor={Ikailo John Sessford},pdfcreator={pandoc + XeTeX},bookmarksnumbered=true}
\pdfstringdefDisableCommands{\def\typesetchar#1{\Uchar#1 }}

% Paragraphs: indented, no extra space, ragged bottom so columns do not stretch.
\setlength{\parskip}{0pt}
\raggedbottom
\emergencystretch=1.5em
\tolerance=1000
\hbadness=10000
\clubpenalty=10000
\widowpenalty=10000

% Section headings. Numbered only when the paper's own headings carry no numbers (build_paper.py decides).
\setcounter{secnumdepth}{-\maxdimen}
\titleformat*{\section}{\normalfont\large\bfseries\raggedright}
\titleformat*{\subsection}{\normalfont\normalsize\bfseries\raggedright}
\titleformat*{\subsubsection}{\normalfont\normalsize\bfseries\itshape\raggedright}
\titleformat*{\paragraph}{\normalfont\normalsize\itshape}
\titlespacing*{\section}{0pt}{2.0ex plus 1ex minus .2ex}{1.0ex plus .2ex}
\titlespacing*{\subsection}{0pt}{1.6ex plus 1ex minus .2ex}{0.6ex plus .2ex}
\titlespacing*{\subsubsection}{0pt}{1.4ex plus 1ex minus .2ex}{0.5ex plus .2ex}
\titlespacing*{\paragraph}{0pt}{1.0ex plus .5ex}{0.6em}

% Code: monospaced at the floor size, wrapped at the column edge and never scaled.
\RecustomVerbatimEnvironment{verbatim}{Verbatim}{fontsize=\footnotesize,breaklines=true,breakanywhere=true,breaksymbolleft={},
  xleftmargin=0.5em,frame=leftline,framerule=0.4pt,rulecolor=\color{Gray},framesep=0.5em}

% Tables: small type, tight columns. Long tables break across columns (xtabular); wide ones float full width.
\setlength{\tabcolsep}{4pt}
\renewcommand{\arraystretch}{1.1}
\newcolumntype{L}[1]{>{\RaggedRight\arraybackslash}p{#1}}
\newcolumntype{R}[1]{>{\RaggedLeft\arraybackslash}p{#1}}
\newcolumntype{C}[1]{>{\Centering\arraybackslash}p{#1}}
% Scaling has a floor: nothing is set smaller than \typesetfloor.
% DERIVED: 8 pt is \footnotesize of this 10 pt class, the smallest running text the template sets
% (tables, code), so a scaled display is never smaller than the smallest text already on the page.
\newdimen\typesetfloor \typesetfloor=8pt
\newsavebox{\typesetfitbox}
\newif\iftypesetfitok
\makeatletter
% \typeset@try{width}{base size}: does \typesetfitbox fit in width, scaled no smaller than the floor?
% Sets \iftypesetfitok, and logs the scale it would apply ("typeset-scale: <scale x 1000>").
\def\typeset@try#1#2{%
  \ifdim\wd\typesetfitbox>#1\relax
    \edef\typeset@scale{\the\numexpr\number\dimexpr#1\relax*1000/\number\wd\typesetfitbox\relax}%
    % the effective size, base size x scale, against the floor
    \ifdim\dimexpr#2*\typeset@scale/1000\relax<\typesetfloor
      \typesetfitokfalse
    \else
      \typesetfitoktrue\typeout{typeset-scale: \typeset@scale}%
    \fi
  \else
    \typesetfitoktrue
  \fi}
% \typesetfit{width}{content}: the content at natural size, scaled down only when wider than width,
% and never below the floor; content too wide even at the floor is set at the floor and overflows,
% which the build reports ("typeset-overflow").
\newcommand{\typesetfit}[2]{\sbox{\typesetfitbox}{#2}\typeset@place{#1}{\f@size pt}}
\def\typeset@place#1#2{%
  \typeset@try{#1}{#2}%
  \iftypesetfitok
    \ifdim\wd\typesetfitbox>#1\relax\resizebox{#1}{!}{\usebox{\typesetfitbox}}\else\usebox{\typesetfitbox}\fi
  \else
    \edef\typeset@floorscale{\the\numexpr\number\typesetfloor*1000/\number\dimexpr#2\relax\relax}%
    \ifnum\typeset@floorscale<1000
      \typeout{typeset-scale: \typeset@floorscale}%
      \scalebox{\strip@pt\dimexpr\typeset@floorscale pt/1000\relax}{\usebox{\typesetfitbox}}%
    \else
      \usebox{\typesetfitbox}%
    \fi
    \typeout{typeset-overflow: \the\wd\typesetfitbox\space into \the\dimexpr#1\relax}%
  \fi}
% \typesetdisplay{math}{split}: a display equation, and the same equation broken into lines at its
% top-level \qquad separators (empty when it has none). The first of these that fits, never below the
% floor, is used: the equation in the column; the split form in the column; the equation across both
% columns (a strip at this point of the text); the split form across both columns.
\def\typeset@show{%
  \[\ifdim\wd\typesetfitbox>\linewidth\resizebox{\linewidth}{!}{\usebox{\typesetfitbox}}\else\usebox{\typesetfitbox}\fi\]}
\def\typeset@strip{%
  \typeout{typeset-strip: display}%
  \par\begin{strip}\vspace{-\abovedisplayskip}\[\typeset@place{\textwidth}{\f@size pt}\]\vspace{-\belowdisplayskip}\end{strip}%
  \noindent\ignorespaces}
\long\def\typesetdisplay#1#2{%
  \sbox{\typesetfitbox}{\(\displaystyle #1\)}%
  \typeset@try{\linewidth}{\f@size pt}%
  \iftypesetfitok\typeset@show\else
    \long\def\typeset@alt{#2}%
    \ifx\typeset@alt\@empty\typeset@strip\else
      \sbox{\typesetfitbox}{\(\displaystyle #2\)}%
      \typeset@try{\linewidth}{\f@size pt}%
      \iftypesetfitok\typeset@show\else
        \sbox{\typesetfitbox}{\(\displaystyle #1\)}%
        \typeset@try{\textwidth}{\f@size pt}%
        \iftypesetfitok\else\sbox{\typesetfitbox}{\(\displaystyle #2\)}\fi
        \typeset@strip
      \fi
    \fi
  \fi}
% \typesetcolumntable{tabular}: an in-column table at natural size, in the column if it fits there
% (tables are already at the floor size, so they are never scaled), otherwise across both columns.
\long\def\typesetcolumntable#1{%
  \sbox{\typesetfitbox}{\footnotesize #1}%
  \ifdim\wd\typesetfitbox>\linewidth
    \typeout{typeset-strip: table}%
    \begin{strip}\centering\typeset@place{\textwidth}{\typesetfloor}\par\end{strip}%
  \else
    \begin{center}\usebox{\typesetfitbox}\end{center}%
  \fi}
\makeatother

% For the build's report: how full each column is when it is made ("typeset-column: page columns
% text floats height", lengths in sp). Float-only pages are made elsewhere and are not logged.
\usepackage{etoolbox}
\makeatletter
\def\typeset@colfill{%
  \setbox\z@\vbox{\unvcopy\@cclv}%
  \typeout{typeset-column: \number\c@page\space\if@twocolumn 2\else 1\fi\space
    \number\dimexpr\ht\z@+\dp\z@+\ht\footins+\dp\footins\relax\space
    \number\dimexpr\textheight-\@colroom\relax\space\number\textheight}}
\pretocmd{\@makecol}{\typeset@colfill}{}{\typeout{typeset-column: unavailable}}
% A heading above the paper's section level (a part of the paper) is set as a centred divider.
\newcounter{typesetpart}
\newcommand{\typesetpart}[2]{\par\vspace{2.5ex plus 1ex}\stepcounter{typesetpart}%
  \pdfbookmark[0]{#2}{typesetpart.\arabic{typesetpart}}%
  {\centering\large\bfseries #1\par}\vspace{1.5ex plus .5ex}}
\makeatother

% Floats: allow full-width floats to take most of a page, and many to queue.
\renewcommand{\topfraction}{0.9}
\renewcommand{\bottomfraction}{0.8}
\renewcommand{\textfraction}{0.07}
\renewcommand{\floatpagefraction}{0.75}
\renewcommand{\dbltopfraction}{0.9}
\renewcommand{\dblfloatpagefraction}{0.75}
\setcounter{topnumber}{3}
\setcounter{dbltopnumber}{3}
\setcounter{totalnumber}{5}
\extrafloats{200}
\setkeys{Gin}{width=\linewidth,height=0.8\textheight,keepaspectratio}

% Pandoc helpers.
\providecommand{\tightlist}{\setlength{\itemsep}{0pt}\setlength{\parskip}{0pt}}
\providecommand{\pandocbounded}[1]{#1}
\newcommand{\passthrough}[1]{#1}

\pagestyle{fancy}
\fancyhf{}
\renewcommand{\headrulewidth}{0pt}
\fancyfoot[C]{\small\thepage}
\fancypagestyle{plain}{\fancyhf{}\renewcommand{\headrulewidth}{0pt}\fancyfoot[C]{\small\thepage}}

\AtBeginDocument{%
\typesetmap{"2192}% →
\typesetmap{"2212}% −
}

\begin{document}

% Title block across both columns, with the abstract under it when the two leave room for the body to
% start on page 1; otherwise the abstract opens the two-column body, so no page is left near-empty.
\makeatletter
\newsavebox{\typesettitle}
\newsavebox{\typesetfront}
\setbox\typesettitle=\vbox{%
  \hsize=\textwidth \linewidth=\textwidth \columnwidth=\textwidth
  \begingroup\centering
  {\LARGE\bfseries Selection for language-model context, read from the data\par}
  \vspace{0.8ex}{\large Using entroptics to decide how many retrieved items to keep, which to keep, and where foreign material sits\par}
  \vspace{1.8ex}
  {\normalsize Ikailo John Sessford\par}
  \vspace{0.3ex}
  {\small\itshape Ikailo Inc.\par}
  {\small \href{mailto:john@ikailo.com}{\texttt{john@ikailo.com}}\par}
  \vspace{0.6ex}{\small Version 0.3.0. November 2026.\par}
  \endgroup
}
\long\def\typesetrule{\vspace{1.6ex}\noindent\rule{\linewidth}{0.4pt}\par\vspace{1.5ex}}
\long\def\typesetabstract{%
A language model's context is assembled from several sources at once: a retrieval-augmented
prompt, an agent scratchpad, a long conversation. The decisions made about it --- how many
retrieved items to keep, which of them to trust, whether an answer stayed on its evidence --- are
conventionally made with one number per item against one summary of the whole: a cosine, a
centroid, a fixed top-\(k\).\\
This work reads the context as the sources it spans, each one the directions it resolves above
its own noise floor, using the \texttt{entroptics} instrument, and makes each of those decisions from
those spans. The number of sources and the directions of each are counted from the data, so
nothing is trained, no threshold is set, and no constant is fitted per corpus.\\
\textbf{Where relevance spreads across several facets, the derived cut holds \(F_1\) at 1.000 through
four facets and at 0.915 and 0.711 at five and six, where a cosine cut falls to 0.000.} On
locating a real prompt injection planted in real documents, it leads both a centroid and a
k-nearest-neighbour detector at every source count measured, by a margin over k-NN that grows
from \(+0.005\) at one source to \(+0.024\) at five, and at five sources it names the attack as the
first item to inspect almost six times as often as a centroid does. On deciding whether a
continuation belongs to its context, it leads a centroid by a margin that grows from \(+0.056\) at
two sources to \(+0.172\) at five, ties k-NN from two sources up, and trails it by \(0.007\) at one.\\
One statement covers all of it. \textbf{These reads beat a single summary statistic when the deciding
distribution is multi-modal, and over the source counts measured the margin grows with how
multi-modal it is.} Where the deciding distribution has one mode they reduce to the statistic
they are measured against and tie; on four BEIR corpora and an end-to-end HotpotQA run they tie
with cosine.
}
\setbox\typesetfront=\vbox{%
  \hsize=\textwidth \linewidth=\textwidth \columnwidth=\textwidth
  \unvcopy\typesettitle
  \vspace{2ex}
  {\centering\small\bfseries Abstract\par}
  \vspace{0.6ex}
  \begin{list}{}{\setlength{\leftmargin}{0.06\textwidth}\setlength{\rightmargin}{0.06\textwidth}%
    \setlength{\listparindent}{1.2em}\setlength{\itemindent}{0pt}\setlength{\parsep}{0.5ex}\setlength{\topsep}{0pt}}
  \small\item\relax\typesetabstract
  \end{list}
  \typesetrule
}
% CHOSEN: the title block and abstract go across both columns only if they leave ten lines of each
% column for the body on page 1. Ten lines is about a sixth of the page: enough that page 1 reads as
% the start of the paper, small enough that only a long abstract moves into the columns.
\ifdim\dimexpr\ht\typesetfront+\dp\typesetfront+10\baselineskip\relax<\textheight
  \twocolumn[\unvbox\typesetfront]
\else
  \twocolumn[\unvbox\typesettitle\typesetrule]
  {\centering\small\bfseries Abstract\par}\vspace{0.6ex}
  {\small\setlength{\parskip}{0.5ex}\noindent\typesetabstract\par}
  \typesetrule
\fi
\makeatother

\section{1. The claim}\label{1-the-claim}

\begin{quote}
\textbf{A read taken over the sources a context resolves decides better than one taken against a
single summary of it when the deciding distribution is multi-modal, and over the source counts
measured the advantage grows with how multi-modal it is.}
\end{quote}

Three measurements carry it, on three different axes of modality: facets of a query's relevance
(§4), sources in a context (§5), and topics spanned by a conversation (§6). On two real-corpus
measurements the reads tie with cosine (§7).

\section{2. A span in place of a point}\label{2-a-span-in-place-of-a-point}

Take a context \(C = \{x_1, \dots, x_n\}\) of direction-normalized embeddings drawn from \(M\)
distinct sources. The centroid

\typesetdisplay{\bar{x} = \frac{1}{n}\sum_i x_i \Big/ \Big\lVert \frac{1}{n}\sum_i x_i \Big\rVert}{}

is a single point on the sphere, and ranking by cosine to it is identical to ranking by average
similarity to the rest of the set --- the standard summary an item is usually scored against.
Where the sources occupy \(M\) well-separated directions, \(\bar{x}\) is their weighted mean and for
\(M > 1\) it points into the gap between them. The cosine of a legitimate item from source \(m\)
then falls roughly as \(\cos\theta\), with \(\theta\) that item's angle from the consensus
direction --- a quantity set by \textbf{how spread the context is}, carrying no information about the
item.

A subspace read puts a span where that point was. Let \(X \in \mathbb{R}^{n \times d}\) stack the
context, let \(X = U\Sigma V^\top\), and take

\typesetdisplay{V_K = [v_1, \dots, v_K], \qquad K = \#\{\,\sigma_j : \sigma_j \text{ resolves above the noise floor}\,\}.}{\begin{gathered}V_K = [v_1, \dots, v_K],\\ K = \#\{\,\sigma_j : \sigma_j \text{ resolves above the noise floor}\,\}.\end{gathered}}

Score an item by the energy it retains under projection:

\typesetdisplay{\phi_{\text{in}}(y) = \frac{\lVert V_K^\top y \rVert^2}{\lVert y \rVert^2} \in [0, 1].}{}

Each resolved source direction lies in \(\operatorname{span}(V_K)\), so a legitimate item from
\emph{any} source scores high: its own direction is represented whether or not the others are.

One span over every source still asks a looser question than the one posed. It scores an item by
how much of it lies in the span of all the sources together, so an off-topic item that falls
between them reads as partly inside. A nearest-neighbour read asks the local question, whether
the item is close to some part of the context, and on the drift task of §6 it is ahead of the
single span from two sources up. The read therefore keeps one span per source.

The \(K\) dimensions of \(V_K\) seed the split. With \(U \in \mathbb{R}^{n \times K}\) the items'
orthonormal coordinates in \(\operatorname{span}(V_K)\), column-pivoted QR on \(U^\top\) picks \(K\)
representative items, the polar factor of their rows rotates \(U\) onto them, and each item joins
its largest rotated coordinate {[}6{]}. Nothing is set: the pivots are the items left most
independent of those already chosen. Each group \(s\) then gets its own span \(V^{(s)}\), built the
same way from the group alone, and

\typesetdisplay{\phi_{\text{in}}(y) = \max_s \frac{\lVert V^{(s)\top} y \rVert^2}{\lVert y \rVert^2}.}{}

An item is in the context when it lies in one of its sources. The margin over a single-point
summary widens as the sources multiply, because the point moves further from all of them while
each source keeps its own span.

\section{\texorpdfstring{3. Choosing \(K\) without choosing anything}{3. Choosing K without choosing anything}}\label{3-choosing-k-without-choosing-anything}

\(K\) is the one quantity that could be tuned, and it is read rather than set. It comes from the
singular spectrum via the \texttt{entroptics} instrument {[}1{]}, which counts the modes standing above the
noise floor derived from the block's own shape. For an \(n \times d\) block of pure noise the
singular values follow a Marchenko-Pastur law with a Tracy-Widom edge; modes above that edge are
the ones the block resolves. The count is a property of the data and moves with the context: in
§5 it reads 6.6 modes at one source and 13.6 at five, without being told how many sources there
were.

A constant level is a mode: in the optics in which the library speaks it is the zero-order beam,
physically present and passed by the aperture like any other. For a set of embeddings that
matters directly, because \textbf{the shared mean direction is the topic the set is about}: it is the
leading coherent mode, and it carries signal. The count is read on the library's whitened frame,
where the noise floor is defined, and the whitening centres each channel on its mean, so it
counts the directions about the shared mean. The directions are taken from that same centred
set, so the directions and their count come from one frame, and the mean direction the centring
removes is added back:

\typesetdisplay{V_K = \operatorname{orth}\,[\,\bar{x},\ w_1, \dots, w_{K-1}\,],}{}

with \(w_j\) the right singular vectors of the centred set and \(K - 1\) the engine's count. A set
with nothing above the floor about its mean is its mean direction alone, \(K = 1\). On the clean
topic sets of Appendix A.1, \(K\) is 3.0 at three topics and 5.0 at five, and 1.1 at one, where
the noise about the mean resolves a direction in some draws.

Absence is handled by masking. A row or column that was never measured is dropped before the
decomposition and never reaches the denominator. A channel of measured zeros is an observation
of no power and counts as a mode. The two are distinct inputs and give distinct values of \(K\),
which is what a caller with sparse features needs to know.

\section{4. How many to keep, when relevance has several facets}\label{4-how-many-to-keep-when-relevance-has-several-facets}

The largest margin in this work is on selection. A query whose relevance is spread across \(R\)
independent facets is the multi-modal case at its cleanest, and \(R\) is a controlled variable:
the facets are planted, so the axis is exact rather than estimated.

Both sides run the same parameter-free lock over the same head screen; the only difference is
whether the screen precedes the lock (\texttt{research/\allowbreak{}benchmarks/\allowbreak{}bench\_\allowbreak{}facets.\allowbreak{}py}, 60 items, 12 gold,
16 heads, 40 draws).

\begin{table*}[tp]
\centering\footnotesize
\typesetfit{\textwidth}{%
\begin{tabular}{lllll}
\toprule
facets \(R\) & cosine \(F_1\) & \textbf{entroptics} \(F_1\) & edge & 95\% CI \\
\midrule
1 & 1.000 & 1.000 & \(+0.000\) & {[}\(+0.000\), \(+0.000\){]} \\
2 & 1.000 & 1.000 & \(+0.000\) & {[}\(+0.000\), \(+0.000\){]} \\
3 & 1.000 & 1.000 & \(+0.000\) & {[}\(+0.000\), \(+0.000\){]} \\
4 & 0.806 & \textbf{1.000} & \(\mathbf{+0.194}\) & {[}\(+0.099\), \(+0.301\){]} \\
5 & 0.000 & \textbf{0.915} & \(\mathbf{+0.915}\) & {[}\(+0.866\), \(+0.955\){]} \\
6 & 0.000 & \textbf{0.711} & \(\mathbf{+0.711}\) & {[}\(+0.624\), \(+0.793\){]} \\
\bottomrule
\end{tabular}}
\end{table*}

\textbf{Entroptics holds \(F_1\) at 1.000 through four facets and keeps 0.915 and 0.711 at five and
six.} The cosine cut matches it while the relevance has one to three facets and falls to 0.000
by five, so the edge opens at four facets and is largest at five. The reason is §2's, transposed
from a context to a query: one scalar carries one facet's worth of evidence, and each gold
item's total weakens toward the distractor level as the facets multiply, where an \(R\)-mode
subspace holds all of them at once.

The count kept is derived per query throughout. Nothing is swept, and the same lock produces
both columns.

\section{5. Which item is foreign, when a context spans several sources}\label{5-which-item-is-foreign-when-a-context-spans-several-sources}

The second axis of modality is the number of sources a context is assembled from. Legitimate
content is whole wikitext-103 articles, so each source is a genuine topic; the attacks are the
\texttt{deepset/\allowbreak{}prompt-{}\allowbreak{}injections} corpus, strings written to attack real systems. One attack is
appended to a context built from \(M\) sources and the read ranks it. Encoder is
\texttt{bge-{}\allowbreak{}small-{}\allowbreak{}en-{}\allowbreak{}v1.\allowbreak{}5}, 200 trials per row, intervals paired per trial
(\texttt{research/\allowbreak{}benchmarks/\allowbreak{}bench\_\allowbreak{}injection.\allowbreak{}py}).

Two baselines. The \textbf{centroid} is the standard single summary of §2. \textbf{k-NN} scores an item by
its similarity to its five nearest neighbours --- a local read, so it carries no penalty for a
context that spans topics, and it is the stronger of the two.

\begin{table*}[tp]
\centering\footnotesize
\typesetfit{\textwidth}{%
\begin{tabular}{llllllll}
\toprule
sources & modes & items & \textbf{entroptics} & centroid & k-NN & vs centroid & vs k-NN \\
\midrule
1 & 6.6 & 39 & \textbf{0.988} & 0.983 & 0.984 & \(+0.005\) {[}\(+0.000\), \(+0.010\){]} & \(\mathbf{+0.005}\) {[}\(+0.001\), \(+0.008\){]} \\
2 & 8.7 & 77 & \textbf{0.972} & 0.902 & 0.959 & \(+0.070\) {[}\(+0.055\), \(+0.087\){]} & \(\mathbf{+0.013}\) {[}\(+0.008\), \(+0.019\){]} \\
3 & 10.6 & 116 & \textbf{0.967} & 0.833 & 0.947 & \(+0.134\) {[}\(+0.111\), \(+0.157\){]} & \(\mathbf{+0.019}\) {[}\(+0.013\), \(+0.026\){]} \\
5 & 13.6 & 191 & \textbf{0.941} & 0.718 & 0.917 & \(+0.223\) {[}\(+0.192\), \(+0.254\){]} & \(\mathbf{+0.024}\) {[}\(+0.016\), \(+0.032\){]} \\
\bottomrule
\end{tabular}}
\end{table*}

\textbf{Against the standard summary the margin grows from \(+0.005\) to \(+0.223\)}, as §2 predicts: the
centroid falls from 0.983 to 0.718 as the sources multiply while entroptics holds from 0.988 to
0.941. At one source the interval reaches zero.

\textbf{Against the local detector the margin grows too, from \(+0.005\) to \(+0.024\)}, every interval
excluding zero. k-NN carries no penalty for a context that spans topics, and the per-source read still gains
on it as the sources multiply.

The operator's question is which item to open first, and there entroptics has the highest rate
at every source count:

\typesetcolumntable{%
\begin{tabular}{llll}
\toprule
sources & \textbf{entroptics} & centroid & k-NN \\
\midrule
1 & \textbf{0.785} & 0.735 & 0.720 \\
2 & \textbf{0.450} & 0.250 & 0.365 \\
3 & \textbf{0.410} & 0.145 & 0.300 \\
5 & \textbf{0.175} & 0.030 & 0.120 \\
\bottomrule
\end{tabular}}

At five sources it names the attack first almost six times as often as the centroid does, from a
pool of 191 items. Each item is scored against every source with itself removed from its own;
Appendix A.1 gives the reason.

\section{6. Whether a continuation belongs to its context}\label{6-whether-a-continuation-belongs-to-its-context}

The third axis is a conversation's topic spread, and the question is non-adversarial: \textbf{is this
continuation on-topic?} Held-out sentences from the context's own sources --- on-topic, and held
out of the context --- are scored against sentences from a different article. Same encoder, 200
trials per row, six held-out sentences per source so the target count resolves every effect
reported (\texttt{research/\allowbreak{}benchmarks/\allowbreak{}bench\_\allowbreak{}drift.\allowbreak{}py}).

\begin{table*}[tp]
\centering\footnotesize
\typesetfit{\textwidth}{%
\begin{tabular}{llllllll}
\toprule
sources & modes & context & \textbf{entroptics} & centroid & k-NN & vs centroid & vs k-NN \\
\midrule
1 & 3.3 & 12 & 0.959 & 0.960 & \textbf{0.966} & \(-0.002\) {[}\(-0.007\), \(+0.004\){]} & \(-0.007\) {[}\(-0.012\), \(-0.003\){]} \\
2 & 4.1 & 24 & 0.946 & 0.890 & \textbf{0.948} & \(\mathbf{+0.056}\) {[}\(+0.046\), \(+0.067\){]} & \(-0.002\) {[}\(-0.007\), \(+0.003\){]} \\
3 & 5.1 & 36 & \textbf{0.937} & 0.827 & 0.936 & \(\mathbf{+0.110}\) {[}\(+0.096\), \(+0.123\){]} & \(+0.001\) {[}\(-0.003\), \(+0.004\){]} \\
5 & 7.0 & 60 & \textbf{0.910} & 0.738 & 0.909 & \(\mathbf{+0.172}\) {[}\(+0.158\), \(+0.187\){]} & \(+0.001\) {[}\(-0.003\), \(+0.005\){]} \\
\bottomrule
\end{tabular}}
\end{table*}

\textbf{The same growth, from an independent task.} The edge over the standard summary opens at two
sources and grows to \(+0.172\) at five, every interval from two upward excluding zero, with the
centroid falling 0.960 → 0.738 while entroptics holds 0.959 → 0.910. §5 appends foreign text and
this withholds native text; they are measured separately and agree on how the margin grows
with the sources.

\textbf{Against the local neighbour read it ties from two sources up}, every interval containing
zero, and trails it at one source by \(0.007\) {[}\(0.003\), \(0.012\){]}. A single span over every
source, the read of §2's first half (\texttt{bench\_\allowbreak{}drift.\allowbreak{}py} reports it), scores 0.964, 0.937, 0.918 and
0.882: it ties k-NN at one source, \(-0.002\) {[}\(-0.005\), \(+0.002\){]}, and trails it by \(0.011\)
{[}\(0.007\), \(0.015\){]}, \(0.018\) {[}\(0.014\), \(0.023\){]} and \(0.027\) {[}\(0.023\), \(0.032\){]} at two, three and
five. The per-source read closes that gap from two sources up and gives up \(0.005\) to the single
span at one source, where the split finds 3.3 groups in one article on average.

\section{7. Real corpora, where relevance has one mode}\label{7-real-corpora-where-relevance-has-one-mode}

The law predicts a tie wherever the deciding distribution has one mode, because the read
reduces to the statistic it is measured against. Two real-corpus measurements confirm it.

\textbf{Four BEIR corpora} --- nfcorpus, scifact, arguana, fiqa, 2,421 scored queries, encoded with
\texttt{all-{}\allowbreak{}MiniLM-{}\allowbreak{}L6-{}\allowbreak{}v2} --- against the identical
\texttt{top\_\allowbreak{}break} lock run on plain cosine salience, so the screen is the only difference
(\texttt{research/\allowbreak{}benchmarks/\allowbreak{}bench\_\allowbreak{}beir.\allowbreak{}py}):

\begin{table*}[tp]
\centering\footnotesize
\typesetfit{\textwidth}{%
\begin{tabular}{lllll}
\toprule
dataset & cosine + lock & \textbf{entroptics} & edge & 95\% CI \\
\midrule
nfcorpus & 0.295 & 0.291 & \(-0.005\) & {[}\(-0.018\), \(+0.009\){]} \\
scifact & 0.571 & 0.558 & \(-0.013\) & {[}\(-0.029\), \(+0.002\){]} \\
arguana & 0.257 & 0.256 & \(-0.001\) & {[}\(-0.007\), \(+0.005\){]} \\
fiqa & 0.344 & 0.341 & \(-0.003\) & {[}\(-0.012\), \(+0.007\){]} \\
\bottomrule
\end{tabular}}
\end{table*}

Four ties, every interval containing zero. At the point the cut sees them, 198 of the 2,421
scored queries (8.2\%) carry four or more relevant documents in their horizon: 144 of
nfcorpus's 260, 4 of scifact's 270, none of arguana's 1,357 and 50 of fiqa's 534. On three of
the four corpora the deciding distribution has one mode and §4's condition is absent.

\textbf{HotpotQA end-to-end}, 150 questions, a language model answering from the selected context
(\texttt{research/\allowbreak{}benchmarks/\allowbreak{}bench\_\allowbreak{}hotpot\_\allowbreak{}e2e.\allowbreak{}py}). The decisive arm is cosine's top-\(m\), where \(m\) is
the count entroptics derived \emph{for that question}, so the only difference is which items were
chosen:

\begin{table*}[tp]
\centering\footnotesize
\typesetfit{\textwidth}{%
\begin{tabular}{lllll}
\toprule
model & fixed \(k=5\) & cosine top-\(m\) & \textbf{entroptics} & entroptics − top-\(m\) \\
\midrule
Qwen2.5-1.5B & 0.490 & 0.432 & 0.432 & \(+0.000\) {[}\(+0.000\), \(+0.000\){]} \\
\bottomrule
\end{tabular}}
\end{table*}

The selection matches cosine's, as the law requires on unimodal relevance.

\textbf{What the derived count buys where the reads tie is the count itself.} The best single
constant \(k\) over the four corpora, top-1, gives up as much as \(0.130\) \(F_1\) against a
corpus's own best fixed \(k\) (nfcorpus, whose best is top-10); reading the count per query
gives up at most \(0.020\). That is the value where there is no sweep to run --- a
service answering against more than one corpus, or a corpus whose content moves.

\textbf{Size the budget to the task.} On HotpotQA the derived count spends 38\% of a fixed top-5's
context tokens, 224 against 593, keeping 2.1 paragraphs where a multi-hop answer needs two gold
paragraphs. Against that fixed top-5 the difference is \(-0.058\) {[}\(-0.123\), \(+0.006\){]}, and what
there is of it is the count rather than the ranking: the matched-budget arm above shows the
ranking already matches the baseline. Where a task needs several passages, set a floor on the count.

\section{8. What this costs}\label{8-what-this-costs}

One thin SVD of the context block, \(O(n d \min(n,d))\), on embeddings the pipeline has already
computed for retrieval. No additional model call, no second encoder, no index. Scoring an item
against the basis is a \(d \times K\) matrix-vector product. For the context sizes here --- tens to
a few hundred items --- the read is well under the cost of the retrieval that produced them.

The split adds one \(n \times K\) decomposition and a \(K\)-step pivoted QR, and each source's span
one decomposition of that source's rows. The anomaly read of §5 rebuilds an item's own source
without it, so it takes one noise-floor read and one decomposition per item, each on a single
source's rows rather than the whole context. Appendix A.1 gives the reason.

\section{9. Scope}\label{9-scope}

\textbf{Coherence reads topic.} An answer that stays on its sources' subject scores high whether or
not it is correct: on the 88 HotpotQA questions whose context kept two or more paragraphs,
right answers scored 0.432 and wrong ones 0.420, a separation of \(+0.011\) with no interval reported. It is a monitor for off-source and lost-thread
behaviour.

\textbf{Scores and rankings.} The reads order items within one context. ``Is there an anomaly in this
set at all'' is a question about a reference distribution, which a caller supplies;
\texttt{anomaly\_\allowbreak{}z} takes one.

\textbf{Measured on two encoders.} The BEIR table is \texttt{all-{}\allowbreak{}MiniLM-{}\allowbreak{}L6-{}\allowbreak{}v2}, every other embedding
number is \texttt{bge-{}\allowbreak{}small-{}\allowbreak{}en-{}\allowbreak{}v1.\allowbreak{}5}, and the end-to-end table is Qwen2.5-1.5B. The mechanism in §2 is a statement about
direction-normalized embeddings generally; the magnitudes are measured on these models, and \(K\)
depends on how a model spreads a topic.

\section{10. Reproducing}\label{10-reproducing}

\begin{verbatim}
pip install -e ".[bench]"
python -m research.benchmarks.bench_facets       # §4
python -m research.benchmarks.bench_injection    # §5
python -m research.benchmarks.bench_drift        # §6
python -m research.benchmarks.bench_beir --datasets nfcorpus scifact arguana fiqa   # §7
\end{verbatim}

The end-to-end table in §7 needs a language model as well:

\begin{verbatim}
pip install -e ".[bench,bench-llm]"
python -m research.benchmarks.bench_hotpot_e2e --n 150
\end{verbatim}

\texttt{bench\_\allowbreak{}facets} runs on CPU in seconds and needs no corpus. The others download their data from
Hugging Face. Every number in this paper carries the conditions it was measured under.

\section{Appendix A. Two choices in the API's shape}\label{appendix-a-two-choices-in-the-apis-shape}

\textbf{A.1 --- An item is never scored against a mode it supplied.} \texttt{anomaly} reads each item against
every source of the set, with its own source rebuilt without it; the other sources never
contained it, and an item that forms a source of its own is read against the others alone.
Scoring an item inside a span it helped build loses on §5's benchmark. The same sources with no
item removed from its own (\texttt{bench\_\allowbreak{}injection.\allowbreak{}py} reports this arm) score AUROC 0.456, 0.461,
0.452 and 0.587 at one, two, three and five sources, trailing k-NN by \(0.528\) {[}\(0.469\), \(0.586\){]}
at one source and \(0.330\) {[}\(0.277\), \(0.384\){]} at five.

What the read scores on clean synthetic sets: one injected unit item with amplitude \(f\) on a
direction no topic uses and \(\sqrt{1 - f^2}\) on one topic; eight items per topic, noise norm
0.25, \(d = 192\), 15 draws, top-1 localization (\texttt{research/\allowbreak{}benchmarks/\allowbreak{}bench\_\allowbreak{}leave\_\allowbreak{}one\_\allowbreak{}out.\allowbreak{}py}):

\typesetcolumntable{%
\begin{tabular}{llll}
\toprule
foreignness \(f\) & 1 topic & 3 topics & 5 topics \\
\midrule
0.7 & 1.00 & 1.00 & 1.00 \\
0.9 & 1.00 & 1.00 & 1.00 \\
1.0 & 1.00 & 1.00 & 1.00 \\
\bottomrule
\end{tabular}}

On the set without the injected item \(K\) is 1.1, 3.0 and 5.0 at one, three and five topics, with
as many sources. The injected item reads 0.49 to 0.50, 0.81 and 1.00 outside every source at \(f\)
= 0.7, 0.9 and 1.0, and the highest other item reads 0.08 in every cell, so the injected item
clears the rest in every draw.

\textbf{A.2 --- A subspace overlap is read against its null.} In high dimension two subspaces that are
not specifically aligned are near-orthogonal, so a raw \texttt{subspace\_\allowbreak{}overlap} reports ``most of the
structure lies outside'' for a random matrix as readily as for a real one, measuring the ambient
dimension. The analytic null is \(k_b/d\) --- a random \(k_b\)-dimensional subspace covers that
fraction of anything --- and \texttt{chance\_\allowbreak{}overlap} returns it. Report the ratio of measured overlap to
chance.

\section{References}\label{references}

{[}1{]} \textbf{entroptics} --- the spectral instrument. Resolved-mode counting against a
Marchenk\allowbreak{}o-\allowbreak{}Pastur/\allowbreak{}Tracy-\allowbreak{}Widom noise floor, and the occupancy read used here to set \(K\).
Apache-2.0. \href{https://pypi.org/project/entroptics/}{https:\allowbreak{}/\allowbreak{}/\allowbreak{}pypi.\allowbreak{}org/\allowbreak{}project/\allowbreak{}entropti\allowbreak{}cs/}

{[}2{]} \textbf{wikitext-103} --- Merity et al., \emph{Pointer Sentinel Mixture Models}, ICLR 2017.

{[}3{]} \textbf{deepset/\allowbreak{}prompt-\allowbreak{}injectio\allowbreak{}ns} --- a corpus of real prompt-injection strings.
\href{https://huggingface.co/datasets/deepset/prompt-injections}{https:\allowbreak{}/\allowbreak{}/\allowbreak{}huggingf\allowbreak{}ace.\allowbreak{}co/\allowbreak{}datasets\allowbreak{}/\allowbreak{}deepset/\allowbreak{}prompt-\allowbreak{}injectio\allowbreak{}ns}

{[}4{]} \textbf{BEIR} --- Thakur et al., \emph{BEIR: A Heterogenous Benchmark for Zero-shot Evaluation of
Information Retrieval Models}, NeurIPS 2021.

{[}5{]} \textbf{HotpotQA} --- Yang et al., \emph{HotpotQA: A Dataset for Diverse, Explainable Multi-hop Question
Answering}, EMNLP 2018.

{[}6{]} \textbf{Pivoted-QR spectral clustering} --- Damle, Minden and Ying, \emph{Simple, direct and efficient
multi-way spectral clustering}, Information and Inference 8(1), 2019.

\section{Citation}\label{citation}

\begin{verbatim}
@techreport{ikailo2026entropticscontext,
  title  = {Selection for language-model context, read from the data},
  author = {{Ikailo Inc.}},
  year   = {2026},
  note   = {https://github.com/Agience/entroptics-context}
}
\end{verbatim}

\end{document}
